\section{Pencegahan CSRF: Token, SameSite Cookie, dan Validasi Origin}

\textbf{Cross-Site Request Forgery (CSRF)} membuat pengguna yang sudah login tanpa sadar mengirim request berbahaya --- misalnya, penyerang menyisipkan form tersembunyi di situs lain yang mengirim POST ke aplikasi target. Karena browser mengirim cookie session secara otomatis, server mengira request tersebut sah. Pencegahan utama: \textbf{CSRF token} --- nilai unik dan random yang disisipkan di setiap form dan divalidasi di server \cite{owaspCsrf}.

\begin{contoh}
\textbf{Studi Kasus: Form Login Aman (Integrasi Keamanan Lengkap)}

Form login menerapkan semua praktik keamanan:
\begin{itemize}
  \item CSRF token di form (mencegah CSRF)
  \item Prepared statement untuk query login (mencegah SQL injection)
  \item \code{password\_verify()} untuk verifikasi password hash
  \item \code{session\_regenerate\_id(true)} setelah login (mencegah session fixation)
  \item \code{htmlspecialchars()} pada output (mencegah XSS)
  \item Cookie session: HttpOnly, Secure, SameSite=Strict
\end{itemize}
\end{contoh}

\begin{webcode}[language=PHP, caption={CSRF token: generate, embed, dan validasi}]
<?php
session_start();

// Generate token
if (empty($_SESSION['csrf_token'])) {
  $_SESSION['csrf_token'] = bin2hex(random_bytes(32));
}

// Di form HTML:
// <form method="POST">
//   <input type="hidden" name="csrf"
//     value="<?= $_SESSION['csrf_token'] ?>">
//   ...
// </form>

// Validasi saat proses form
if ($_SERVER['REQUEST_METHOD'] === 'POST') {
  $token = $_POST['csrf'] ?? '';
  if (!hash_equals($_SESSION['csrf_token'], $token)) {
    http_response_code(403);
    die("CSRF token tidak valid.");
  }
  // Token valid, proses form...
  // Regenerate token setelah dipakai
  $_SESSION['csrf_token'] = bin2hex(random_bytes(32));
}
?>
\end{webcode}

